% ===========================================================================
% 04a-ruler.tex -- sec:res-metrics and sec:res-blind (legacy anchors kept).
% ===========================================================================

\beatsection{E}{1}{The metric is the first result}{sec:res-metrics}
\label{sec:methods-metrics}   % legacy anchor, so v1 cross-references resolve

\claimlede{The field's standard score is saturated by a predictor that carries
no information, and of the five metrics audited only \NIR{} grades in the
direction its name implies.}

\begin{question}[1]
Every claim in this chapter is a number handed down by a scoring rule, so the
scoring rule is what has to be audited first. The instrument grades metrics, not
models. It asks each candidate to separate a positive control from a predictor
that knows no drug, and four of the five candidates fail.
\end{question}

\noindent
A reported score is two claims at once, one about the model and one about the
rule that graded it. Nobody checks the second. Two questions decide it. Does the
score reward the thing it is named after, or can it be earned another way? And
does it register a true positive when there is one?

\paragraph{The score the field reports} That score is a correlation between
predicted and true change from control, restricted to the genes that actually
move. The benchmarks of \cref{sec:lit-dispute} report it,
Miller et al.~\cite{miller} find it poorly calibrated, DrEval~\cite{dreval}
counts it among six standing obstacles. Write $x$ for the control profile, $y$
for the treated truth and $\hat y$ for the
prediction,\seealso{\cref{sec:methods-data} fixes $G$.} and let
\begin{equation}
 S_K \;=\; \operatorname*{arg\,top-}K_{\,g} \; \lvert y_g - x_g \rvert
\end{equation}
be the $K$ genes with the largest true change. Then
\begin{equation}
 \DEdr \;=\; \operatorname{corr}\!\big( \left(\hat{y} - x\right)\big|_{S_K}, \;
 \left(y - x\right)\big|_{S_K} \big).
\end{equation}
All three are the panel profiles converted to ranks, absent genes tied at the
worst rank $G$ (\cref{sec:app-convention}), so $S_K$ is selected on rank change
too; the \emph{expression frame} label used elsewhere names the object scored, a
cell profile rather than a residual, and not the units. We report $K = 50$,
Pearson; $K \in \{20, 100, 200\}$ and Spearman move the
numbers by under $0.02$, except at $K = 200$, where the top-$K$ set starts to
include genes that do not move.

\paragraph{What the metric actually rewards} Score a predictor that contains
nothing. Place every gene at the middle rank.\gloss{revert\_center}{Predicts
every panel gene at the panel's middle rank, defined from the metric alone and
needing no model, training set or atlas.} It scores $\DEdr(K{=}50) =$
\ci{0.999913}{0.999908}{0.999919}, above the two-real-cell positive control of
$0.76$ and above the model's $0.73$. Less information yields a higher score,
the hallmark of an artefact; \cref{fig:calibration}, panel a, sets the five arms
against that reference on one axis.

Its partial-$\DEdr$ and $\paneltau$ are undefined, not merely poor, for two
separate reasons. The prediction is the constant midpoint rank $\hat y = G/2$,
so the shift from control is $G/2 - x$, an exact affine function of $x$ from
which nothing survives partialling. And every predicted gene being tied, no rank
ordering exists whose Kendall $\tau$ can be computed. The zero-shift
construction, easily confused with this one, is \texttt{control\_copy},
$\hat y = x$; \texttt{revert\_center} is not it.

\begin{table}[htbp]
\centering
\begin{threeparttable}
\caption{\textbf{$\DEdr$ is saturated by a control artefact.} The two arms
needing no drug information sit above the positive control and above the model;
regressing the control out leaves $\approx 0.26$ of drug-agnostic skill in the
two rows where the partial is defined.}
\label{tab:de-artifact}
\footnotesize
\setlength{\tabcolsep}{4pt}
\begin{tabular}{@{}llrl@{}}
\toprule
\thead{Predictor} & \thead{Information} & \thead{$\DEdr$ (K50)}
  & \thead{partial-$\DEdr$} \\
\midrule
\texttt{revert\_center} (genes $\rightarrow G/2$) & none & \textbf{0.9999}
  & undefined \\
\texttt{revert\_mean} (predict mean control) & none, no fit & 0.961 & 0.25 \\
ridge linear (control $\rightarrow$ shift) & control & 0.947 & 0.28 \\
\addlinespace
positive control (two real cells) & -- & 0.76 & -- \\
model (1B LLM) & control $+$ drug & 0.73 & not measured \\
\bottomrule
\end{tabular}
\begin{tablenotes}[flushleft]\footnotesize
\item The $\DEdr$ column is tier 2; the partial column, tier 1.
\end{tablenotes}
\end{threeparttable}
\end{table}

\paragraph{Two innocent properties compose into the failure} $S_K$ is selected
using the truth, so the exam questions come from the answer key. And treated
cells in this atlas regress toward a shared response, so a prediction reverting
toward the population centre correlates with $y - x$ by construction.
Partial-$\DEdr$ regresses the control out of both arguments first, closing that
route and collapsing the raw $0.95$ to $\approx 0.26$ in the two baselines with
a defined partial value, one of which performs no fitting at all.

The model's own partial-$\DEdr$ was never computed,\footnote{It requires a GPU
pass that was not run, so the cell in \cref{tab:de-artifact} is marked, not
filled with an expected value.} so the comparison anchors to $\paneltau$,
Kendall's $\tau$ over the full $\num{946}$-gene panel, measured for every arm.
The calibration audit finds $\paneltau$ inverted, so neither an absolute claim
nor an ordering read off it counts as evidence of skill. What it can still show
is an \emph{absence} of separation. \texttt{revert\_mean} scores $0.268$, ridge
$0.265$ and the model $0.257$ (tier 2, \texttt{worst} convention), three arms
inside a spread of $0.011$.

The shift therefore decomposes into three parts: control reversion, an
artefact; a generic drug programme worth $\approx 0.26$, trivially matched; and
a drug-specific component no arm is shown capturing. Later residual-frame
instruments show the checkpoint reads trained-drug information while still
failing condition-specific prediction.

\paragraph{One convention, disclosed and then dropped} A cell sentence lists
only expressed genes, so a rank-based metric needs a convention for the unranked
panel genes: \texttt{worst} gives them the last rank $G$, \texttt{francesca} a
fixed middle rank $G/2$. Spike-in $\DEdr$ reads $0.952$ against $0.954$ and
nothing here moves by more than $0.01$, so all figures are quoted under
\texttt{worst}. That diagnostic went to no artifact, so it is a check that was
run, not a reproducible result.

\paragraph{What travels beyond this atlas is the audit, not the verdict} The
saturating predictor is defined from the metric alone, so any result on this
family can be checked the same way. Whether the reversion property is strong
enough to invert the score elsewhere, or in the log-expression units the
benchmarks compute the same functional form in, this cache cannot answer.

\paragraph{The metric this investigation relies on instead} $\NIR$ is Wu
et al.'s ``transposed-rank'' metric from PerturBench~\cite{perturbench}, which
Miller et al.~\cite{miller} identify as one of the few consistently calibrated
metrics across 14 Perturb-seq datasets. It asks not how closely a prediction
matches its own truth, but whether it matches its own truth \emph{more closely
than it matches anyone else's} --- much harder to answer by accident.

\begin{defn}{Normalised Inverse Rank ($\NIR$)}
For a prediction $\hat y$ of condition $i$, with truths $\{y_j\}$ over a declared
comparison set $J$,
\begin{equation}
 \NIR(\hat{y}, i) \;=\; \frac{1}{|J|}\sum_{j \in J}
 \Big[\mathbf{1}\{\sigma(\hat{y}, y_i) > \sigma(\hat{y}, y_j)\}
 + \tfrac{1}{2}\,\mathbf{1}\{\sigma(\hat{y}, y_i) = \sigma(\hat{y}, y_j)\}\Big],
\end{equation}
the fraction of the other conditions in $J$ whose truth the prediction resembles
less than its own, ties counted at one half. $\sigma$ is a \emph{similarity}, so
higher means closer; built from a distance, as with
$\sigma(\hat y, y) = -\lVert \hat y - y \rVert_2$ for \texttt{nir\_expr}, it is
that distance negated. Only $\sigma$ and $J$ change across frames, both
declared where a figure is quoted.
\end{defn}

\Cref{fig:nir} draws one comparison set under this definition, with the three
reading rules that follow from it, the half-tie rule among them.

% fig-nir.tex -- NIR as a discrimination score. Goes with the metric definition
% in Sections/04a-ruler.tex, and is \input from there.
%
% No data dependency: this is a schematic. Its one numeral, NIR = 6/8, is read
% off the drawing itself -- six of the eight hollow circles lie left of the
% filled dot -- and agrees with the definition and the tie-term paragraph in
% 04a-ruler.tex. The sibling fig-nir.json records that check.
%
% This is the v2 COPY of thesis/figs/fig-nir.tex. The v1 file is frozen and was
% not touched. Two changes, and nothing else.
%
% (1) MEASURE. The v1 drawing was 302.79pt wide against v2's 404.03pt (142mm)
%     block -- 75% of the measure, so it sat as a small island with 36mm of
%     white at its right. It is redrawn at the full measure: every x-coordinate
%     is v1's times 1.598, and the marks that ARE the drawing (dot radii, band
%     height, brace amplitude, arrowhead) scale with them, so the figure is
%     enlarged rather than stretched. Type is NOT scaled: \scriptsize and \small
%     stay at the document's own sizes. That is the whole point -- v1's figures
%     were scaled by LaTeX and their type shrank with them. The vertical rhythm
%     of the three reading rules opens from 0.34cm to 0.42cm, proper leading for
%     8pt type rather than the 9.7pt v1 had.
%
%     The explicit bounding box below belongs to this change. A TikZ bounding
%     box includes each node's invisible inner sep, so a picture whose BOX is
%     flush with the measure has its INK sitting 3.5mm inside it at both edges,
%     and the figure reads as not quite reaching the margin. Declaring the box
%     to be the ink fixes that. Measured on the built page: ink 141.82mm inside
%     the 142.00mm block, 0.31pt clear of the left margin and 0.21pt of the
%     right. Nothing is scaled in LaTeX -- no \resizebox, no scale= key, no
%     width= key -- and the document builds with zero overfull boxes.
%
% (2) LEADER. "own truth" floated above the filled dot with a hollow circle
%     1.4cm to its right, and at this width a reader can take it to be labelling
%     that hollow circle. A 0.35cm leader in quietink now ties the label to the
%     dot it names. A plain line, not an arrow: it identifies, it does not
%     assert a direction.
%
% NOT changed: the nine circles and their irregular spacing, NIR = 6/8, all
% three reading rules including the half-tie rule (\S4.7 turns on the half-tie
% property, so its in-figure statement has to survive), and the caption.
\begin{figure}[htbp]
\centering
\begin{tikzpicture}[
  x=1cm, y=1cm, font=\small,
  lbl/.style={font=\scriptsize, text=black!70},
  strong/.style={font=\scriptsize, text=black!90},
  own/.style={circle, fill=black!85, inner sep=3.2pt},
  other/.style={circle, draw=black!55, fill=white, inner sep=3.05pt},
  ar/.style={-{Stealth[length=5.5pt]}, black!65},
  leader/.style={quietink, line width=0.5pt},
]

% The box is the INK, not the ink plus every node's invisible inner sep: left is
% the axis line's left end, right is the right edge of the NIR label's glyphs,
% and the vertical pair reproduces the natural bounding box. 14.190cm in a
% 14.200cm block. Declared FIRST, so nothing drawn below can enlarge it; every
% mark in the picture sits inside it, so nothing spills into the margin either.
\path[use as bounding box] (0.713,-3.080) rectangle (14.903,1.290);

\node[lbl, anchor=south] at (3.28,0.74) {other drugs in the same comparison set};
\node[strong, anchor=south] at (10.07,0.74) {own truth};
% change (2): the leader. 0.35cm, quiet ink, label's south down to the dot's edge.
\draw[leader] (10.07,0.74) -- (10.07,0.39);

\fill[black!10, rounded corners=1.5pt] (0.88,0.00) rectangle (10.07,0.44);
\foreach \x in {1.52, 2.80, 4.15, 5.67, 7.11, 8.71} {\node[other] at (\x,0.22) {};}
\node[own] at (10.07,0.22) {};
\node[other] at (11.43,0.22) {};
\node[other] at (12.54,0.22) {};
\node[anchor=west, font=\small] at (13.325,0.22) {$\NIR = \tfrac{6}{8}$};

\draw[ar] (0.72,-0.34) -- (13.02,-0.34);
\node[lbl, anchor=north] at (6.87,-0.48)
  {similarity of the prediction to each condition's truth $\longrightarrow$};

\draw[decorate, decoration={brace, mirror, amplitude=4.5pt}, black!55]
  (0.88,-1.12) -- (10.07,-1.12);
\node[strong, anchor=north] at (5.43,-1.30) {the shaded fraction \emph{is} the score};

\node[lbl, anchor=west] at (0.72,-1.96) {own truth ranked first $\Rightarrow \NIR \to 1$};
\node[lbl, anchor=west] at (0.72,-2.38) {own truth ranked at random $\Rightarrow \NIR = 0.5$};
\node[strong, anchor=west] at (0.72,-2.80)
  {every similarity \emph{identical} $\Rightarrow \NIR = 0.5$ by the half-tie rule, not $0$};

\end{tikzpicture}%
% The comment character above is load-bearing. Without it the line break is an
% interword space inside the centred line, which spends 3.4pt of a measure the
% drawing now fills exactly and pushes the drawing 1.7pt off centre.
\caption[NIR is a discrimination score]{\textbf{$\NIR$ asks whether a prediction resembles its own
condition's truth more than it resembles the truths of the other conditions in the comparison set
$J$.} Own truth ranked first gives $\NIR \to 1$; ranked at random gives $0.5$; and a predictor whose
similarities are all identical --- the zero vector, which is what ``predict the average drug
response'' becomes in residual space --- scores $0.5$ by the half-tie rule, not $0$.}
\label{fig:nir}
\end{figure}


\paragraph{Why chance is $0.50$, and what that claim rests on} Chance sits at
$0.50$ under a symmetry assumption, not as a consequence of using the same truth
set, which is necessary but insufficient. Three things must hold for an
uninformative prediction: $y_i$ exchangeable with the comparison truths
$\{y_j\}$; its label conferring no privileged relation to the prediction; and
symmetric tie handling. The reference is assumed, not proved, so it is checked
empirically in each frame. In the expression frame the drug-agnostic linear map
and control-copy score $0.500$ and $0.504$; in the residual frame random,
control-copy and \texttt{orth} score $0.501$, $0.506$ and $0.501$. One asymmetry
survives \emph{between} arms. The split-half reference is scored against
half-sample truths, every baseline against full-sample
truths,\seealso{\cref{sec:methods-data} lists every reference in the chapter.}
which makes the comparison conservative.

\paragraph{The tie term is not cosmetic} Discard it as bookkeeping and the
metric breaks in residual space (\cref{sec:methods-residual}), where a predictor
of the generic response is exactly the zero vector, so every comparison is a
tie. A strict inequality scores it $0.000$ instead of the
$0.500$ it deserves, below every random predictor.

\paragraph{The comparison set must be declared too} Conditions are grouped by
(cell line, plate) in the expression frame and by cell line in the residual
frame, and the two are not interchangeable. Widening $J$ across plates moves a
drug-agnostic baseline from $0.161$ to $0.399$ (\cref{sec:res-blind}), more than
double the score for the same predictor. And because the generic
response sits in every truth in equal measure, it cancels, leaving only the
drug-specific difference to decide the score --- the component the reversion
artefact exploits.

\paragraph{Three similarities, one of them graded} Three similarity functions
are computed and stored: \texttt{nir\_expr}, Euclidean distance in expression
space; \texttt{nir\_rank}, the same over rank vectors; and, in the residual
frame of \cref{sec:methods-residual}, cosine
similarity between signed rank-discounted residual vectors. Every
\emph{expression-frame} $\NIR$ figure here is \texttt{nir\_expr}.

\texttt{nir\_rank} never entered the graded family; a cheaper admissibility
screen excludes it.\gloss{admissibility screen}{A pass/fail gate applied before
grading. Can a disjoint sample of a condition's own cells be retrieved above
chance at all?} On the within-plate tier-2 run that later delivers the
drug-blindness verdict ($n = \num{606}$, \texttt{nir\_sameplate.json}), every arm
sits below the $0.50$ chance line, the positive control included at $0.380$,
against control-copy $0.337$, model $0.313$, drug-agnostic linear map $0.285$,
mean baseline $0.070$. That is no argument against rank-based metrics;
$\NIR$'s rank orders \emph{conditions}, not genes.

\paragraph{Grading the metrics, not the models} The framework is Miller
et al.'s~\cite{miller}. Ask not which model scores highest but which
\emph{metric} can see quality at all.

\begin{defn}{Dynamic Range Fraction ($\DRF$)}
\begin{equation}
 \DRF(m) \;=\; \frac{m(\textrm{pos}) - m(\textrm{neg})}%
 {m(\textrm{perfect}) - m(\textrm{neg})},
\end{equation}
on a scale where $0$ is the drug-agnostic baseline and $1$ a perfect predictor.
Here $\textrm{pos}$ is an interpolated-duplicate positive control, a technical
duplicate interpolated toward the mean baseline with per-gene weights from
differential-expression significance, and $\textrm{neg}$ is a drug-agnostic
leave-one-out mean. A negative $\DRF$ means the metric scored the uninformative
baseline above the positive control in that run.
\end{defn}

Reading a negative value as \emph{inversion} takes one further step, that the
sign belongs to the metric and not to the run.\gloss{inversion}{A metric
placing the drug-agnostic baseline above the positive control, as a property of
the metric and not of one run.} Three of the four negatives survive it; the
fourth is carved out where it appears.

A drug-agnostic mean of that kind also appears in the leak audit of
\cref{sec:res-blind}, scoring $0.161$ within plate and $0.399$ across; the
$\DRF$ audit's own negative scores $0.457$ across plates and $0.272$ within.
Whether the two constructions coincide is \needsaudit{under audit}, and nothing
depends on it. $\DRF$ is formed inside one run from that run's own positive and
negative, so all it needs is that the negative carry no drug information, which
both do by design.

\paragraph{Four choices decide what the number means} Three govern the point
estimate. Ties take average ranks, since otherwise a tied gene's rank would
follow its position in the panel. The ratio is recomputed whole inside every bootstrap
resample, since holding a denominator estimated from the same data understates
the interval. And Holm control runs across the five metrics at once.

The fourth governs whether the interval means anything, and in the stored runs
it was not made.\corrmark{The $\DRF$ bootstrap resampled the wrong unit and
imposed no null; both repaired and both arms rerun before any figure here was
read.} The resampler treated each of the all-plates run's $\num{1820}$ rows (one
drug within one cell line) as its own cell line, against the $25$ that exist,
and the same-plate run carried the identical defect. The one-sided $p$ was read
off the observed resampling distribution with no null imposed, so it degenerates
to the floor $1/2001$ where $\DRF > 0$ and to $1$ where $\DRF < 0$; no metric
was in fact tested. Both are repaired and both arms rerun. The rerun
also moved the cluster to the cell line, since clustering on the (cell
line, plate) pair repeats the original error one level up; grouping unit and
cluster unit are reported separately.

\begin{table}[htbp]
\begin{fullwidth}
\begin{threeparttable}
\caption{\textbf{Under the hardened protocol, $\NIR$ is the only one of the five
whose $\DRF$ is positive, and the sign pattern is identical in both arms.} All
four rivals score the drug-agnostic baseline above the positive control and
exclude zero in both arms. The second column is what separates $\NIR$.}
\label{tab:drf}
\small
\begin{tabularx}{\linewidth}{@{}l l r r
  >{\raggedright\arraybackslash}X@{}}
\toprule
\thead{Metric} & \thead{Scored against} & \thead{$\DRF$ all plates}
  & \thead{$\DRF$ same plate} & \thead{Reading} \\
\midrule
$\NIR$ (\texttt{nir\_expr}) & other conditions' truths & $+0.635$ & $+0.545$
  & the only positive \\
\addlinespace
\texttt{weighted\_r2} & its own truth & $-0.163$ & $-0.104$
  & no directional claim \\
$\paneltau$ & its own truth & $-0.319$ & $-0.250$ & inverted \\
\texttt{spearman\_expr} & its own truth & $-0.650$ & $-0.500$ & inverted \\
$\DEdr$ & its own truth & $-0.686$ & $-0.515$ & inverted; the default \\
\bottomrule
\end{tabularx}
\begin{tablenotes}[flushleft]\footnotesize
\item The all-plates run has $\num{1820}$ rows (one drug within one cell line)
clustered on the $25$ cell lines they nest in; the same-plate run has
$\num{586}$ rows over $44$ cell-line$\times$plate groups spanning $27$ cell
lines, clustered on cell line. $\NIR$ reads \ci{+0.635}{+0.604}{+0.663} and
\ci{+0.545}{+0.516}{+0.569}, Holm $p = 0.002$ at the resampling floor in both
arms. The four negatives are point estimates, each excluding zero on the same
clustering in both arms; \texttt{weighted\_r2} sits closest to zero and its sign
tracks cell count, hence no directional claim.
\end{tablenotes}
\end{threeparttable}
\end{fullwidth}
\end{table}

\paragraph{Only one candidate points forward} \Cref{tab:drf} carries the values;
three things it does not. Against a null imposed by centring the resampling
distribution at $\DRF = 0$, no draw of $\num{2000}$ reaches $\NIR$'s observed
value, so its $p$ sits at the resampling floor and the estimate is some forty
cluster-robust standard errors from zero. Restricting swap partners to the same
plate changes no sign and matches the arms at $27$ cell lines within plate
against $25$ across; an earlier configuration capped the run at $25$
\emph{groups} and admitted only $18$ distinct cell lines, and still gives $\NIR$
\ci{+0.519}{+0.479}{+0.546} over $335$ rows. And the inversion belongs to the
data, not to a leak. A leave-one-out correction barely moves the numbers, real
cell lines carrying 20--204 drugs, and it survives holding the plate constant.
What the argument uses is the sign pattern, drawn for the same-plate arm in
\cref{fig:calibration}, panel b; no ranking among the negative metrics is
claimed.

\begin{figure}[htbp]
\centering
% drawn at 142mm by thesis_v2/figs/calibration.py, the width of the measure; no
% width= key, so it is placed at scale 1.0 and its type prints at the point
% sizes it was set in.
\includegraphics{fig-calibration}
\caption[A zero-information predictor outscores the model]{\textbf{A predictor
carrying nothing about drug, cell or biology scores the field's standard metric
above both the within-well reference and the model, and of five metrics only
$\NIR$ grades in the direction its name implies.} (a) $\DEdr$, $K = 50$,
Pearson, tier 2 unseen drugs, \texttt{worst} convention, $n = 400$ conditions;
the axis runs from $0$ and is not truncated. The grey band marks the within-well
split-half precision reference at $0.76$, the two-real-cell positive control of
\cref{tab:de-artifact}: two halves of the same well, a precision figure and not
a biological replicate. Its width is a drawing device, not an interval, and the
arm marks are bare --- the source carries a normal-approximation interval over
conditions, and this document draws only the clustered bootstrap, so nothing is
drawn rather than an interval of another kind. (b) $\DRF$ for the five metrics
of \cref{tab:drf}, the same-plate arm only: $586$ drugs over $44$
cell-line$\times$plate groups spanning $27$ cell lines, the pale block at each
arrowhead the $95\%$ percentile bootstrap interval clustered on cell line
quoted there.}
\label{fig:calibration}
\end{figure}

\paragraph{What this replicates, and where it does not} The methodological core
of Miller et al.\ replicates, and this audit sharpens it. Metric choice
dominates, and the
control-referenced correlation they flag is the worst performer in both arms
($\DEdr$, $-0.686$ across plates and $-0.515$ within). Their calibrated set does
not transfer intact. \texttt{weighted\_r2}, this project's counterpart of the
weighted $R^2\Delta$~\cite{mejia} they endorse, is negative in both audited
arms, though closest to zero and sign-tracking cell count, so it carries no
directional claim; the inversions of $\paneltau$ and \texttt{spearman\_expr}
belong to neither calibrated set. Their endorsed $\NIR$ does survive. Being
rank-based is not what protects it; refusing to score anything the generic
response can supply is.

\paragraph{Adopting the instrument is also a defence} Miller et al.\ conclude
that under calibrated metrics deep models outperform uninformative baselines;
\cref{sec:res-blind} finds ours at chance on that very
metric.\seealso{\cref{sec:lit-dispute} sets out the dispute this audit enters.}
Genetic against drug perturbation and a discriminative against a generative
model are candidate explanations; \cref{sec:res-lookup} supplies a third.
Adopting their metric also puts the verdict beyond the objections their own
paper raises.

\paragraph{The residual frame is audited separately} A metric is not calibrated
by having a relative calibrated, so the residual frame carries its own
$\DRF$.\framecue{\Rframe}{the second audit is scored in the residual frame,
\cref{sec:methods-residual}} With a split-half replicate as positive control and
the evaluation's drug-agnostic arms as negatives, it reads
\ci{+0.704}{+0.667}{+0.736} against \texttt{control\_copy}, clustered on the
$42$ cell lines. Which negative is chosen barely matters. A random vector and
\texttt{orth} both give $+0.707$, and all three sit at chance as a matter of
measurement ($0.506$, $0.501$, $0.501$). The positive control is the more
conservative of the two, a raw split-half replicate instead of the denoised
interpolated duplicate, so the residual frame is the better separated ($+0.704$
against $+0.635$) despite the harder reference. \texttt{generic} scores exactly
$0.500$ with zero variance, the residual being defined by subtracting it, and
carries no weight.

\paragraph{One declared exception on the intervals} $\DRF$ is a ratio of means;
the two-way estimator used everywhere else here is written for a
mean.\bench{two-way cluster-robust intervals on every
mean}{\cref{sec:methods-data}} Both $\DRF$ intervals are therefore
\emph{one-way} percentile cluster bootstraps that resample whole cell lines and
recompute the ratio inside each draw, respecting the cell-line dependence and
ignoring the crossed well dimension, so they are narrower than fully crossed
intervals. The separation they grade is far too large for that to change the
sign, but the caveat belongs beside the number.

\paragraph{Two questions about the instrument remain, and they differ} The first
is whether the conditions are empirically discriminable at the tested budget. A
spike-in forced choice answers it. Within a plate, a reference pseudobulk from
drug A's cells must be scored closer to a disjoint sample of A's own cells than
to drug B's, chance $0.50$, with contamination of B by A titrating the
difficulty. At zero contamination three of the six metrics run on it
separate the two populations essentially perfectly ($\paneltau = 1.000$,
$\DEdr = 0.995$, top-$n$ $\tau = 0.987$, plate held constant); at full
contamination those same three return $0.501$--$0.502$. The whole-profile
similarities are much weaker even at zero contamination, cosine $0.740$,
Spearman $0.680$, cosine on the shift $0.620$, which matters because a
whole-profile similarity is what $\NIR$ uses here. Split-half retrieval
establishes discriminability, not predictability from untreated input, so
discriminability is not the sole bottleneck at this budget.

The spike-in cannot answer the second question, whether $\NIR$ registers a true
positive, since $\NIR$ was not among the metrics run on it. The evidence is the
$\DRF$ audit's own positive control. Under $\NIR$ the interpolated
duplicate scores $0.802$ against the drug-agnostic mean's $0.457$ across plates,
and $0.668$ against $0.272$ within plate; in the residual frame a split-half
replicate reaches $0.854$ against negatives at $0.501$--$0.506$. None of those is
the reference the verdict is measured against. That reference is the within-well
split-half precision estimate recomputed on the tier-2 within-plate benchmark,
where it sits above chance but not far above it; that, and not $1.0$, is the bar
\cref{sec:res-blind} states the model against.

\begin{scopelimit}[carried into \cref{sec:res-blind} and the rest of the chapter]
The audit licenses $\NIR$ and nothing else. $\paneltau$ is inverted in both
audited arms, so no absolute claim rests on it and no ordering read off it counts
as evidence of skill; an inverted metric preserves neither. Where $\DEdr$
appears later it is only as a paired comparison of two arms on one scale, never
as an absolute score. And a metric is not calibrated by having a relative
calibrated. What is graded here is \texttt{nir\_expr} in the expression frame and
the residual frame's cosine similarity in its own frame, each against its own
controls. A third $\sigma$ would need its own audit.
\end{scopelimit}

\begin{claim}
This audit grades five candidates, and $\NIR$ is the only
one with a positive $\DRF$ (\ci{+0.635}{+0.604}{+0.663} on a one-way cluster
bootstrap over the $25$ cell lines, the declared exception to this chapter's
two-way rule, Holm $p = 0.002$; \ci{+0.545}{+0.516}{+0.569} over $27$ cell lines
with same-plate partners, Holm $p = 0.002$), and the other four are negative with
intervals excluding zero in both audited arms. Three of them ($\paneltau$,
\texttt{spearman\_expr} and $\DEdr$) are inverted rather than merely
uninformative, each ranking the drug-agnostic baseline above the positive
control; \texttt{weighted\_r2} sits closest to zero and its sign tracks cell
count, so no directional claim is made for it. The saturating predictor is not a
curiosity. It carries no information about drug, cell or biology and scores
$\DEdr = 0.9999$ in rank space, above the two-real-cell positive control of
$0.76$ and above the model's $0.73$. Under $\NIR$ the positive control is retrieved well clear of
the negative ($0.802$ against $0.457$ across plates, $0.668$ against $0.272$
within), so the instrument registers a true positive. Every drug-use claim in
the rest of this chapter that rests on a continuous prediction score is
therefore made on the calibrated metric; $\DEdr$ appears below only where two
arms are compared with each other on the same scale, never as an absolute score
of the model's drug use. The remaining instruments (forced-choice grading,
output invariance, and the logit-space activation tests of
\cref{sec:res-mechanism}) are deliberately on scales of their own, each with a
fixed chance value or a paired control.
\end{claim}

% ===========================================================================
% Beat 2.
% ===========================================================================

\beatsection{E}{2}{The model is drug-blind}{sec:res-blind}
\label{sec:methods-controls}   % legacy anchor, so v1 cross-references resolve

\claimlede{On the calibrated metric the model is drug-blind in the expression
frame, including on the identifiable stratum.}

\begin{question}[2]
With a ruler that cannot be gamed in hand, the investigation can ask the
question it exists to answer. Does the model use the drug at all? Four
instruments answer it, and each exists because one specific objection would
otherwise be fatal.
\end{question}

\noindent
The answer is no, and the instruments matter only as the objections they close.
On the calibrated metric, within plate on tier 2 ($n = \num{606}$), the model
scores $0.498$ against a within-well split-half precision reference of
$0.576$.\bench{one estimator,
seven settings}{\cref{tab:ceilings}} It is matched by predictors that know
nothing about the drug: a drug-agnostic linear map scores $0.500$, a
control-copy $0.504$, the global mean $0.180$. And on the identifiable stratum,
whose own within-well reference clears the pre-set $0.80$ bar, a
zero-drug-information control-copy reaches $0.766$ against the model's $0.768$.

\paragraph{Four instruments, and each closes a different objection} The
benchmark itself, within plate; the scramble ablation, the only manipulation
that isolates drug use; forced-choice grading, whose chance value is fixed; and
output invariance, which no scoring choice can confound. They are exposed to
different artefacts, so no single one explains all four.

\paragraph{One measurement rule comes first} Drug and plate are partially
confounded by the atlas's design
(\cref{sec:methods-data}),\framecue{\Eframe}{all four instruments are scored in
the expression frame} so grouping by cell line alone lets the batch effect
stand in for drug identity. A drug-agnostic mean baseline, containing no drug
information whatsoever, scores $0.399$ when the comparison set may span plates
against $0.161$ when it may not, and the precision reference falls from $0.625$
to $0.575$ alongside it, the batch
effect having inflated that too.\seealso{\cref{tab:plateleak} is the leak
audit, a wider run than the benchmark} That is plate leakage, measured and not
assumed. Being wider, its within-plate figures sit slightly off the ones above;
what it licenses is the rule, not a substitute set of scores.

\begin{scopelimit}[carried into every discrimination result in this chapter]
Every discrimination instrument gained a \texttt{--same\_plate\_only} mode and
every discrimination result in this chapter was re-run within plate, except
where a departure from that rule is named at the point of use, as with the
swap-distance sweep that answers the fourth objection below, whose swap partners
were not plate-restricted (\cref{sec:app-sweep}), and the residual-frame scores
of \cref{sec:methods-residual}, whose comparison set is grouped by cell line and
whose justification is given there. Conclusions are unchanged; some magnitudes
shrink.
\end{scopelimit}

\paragraph{The central manipulation is the scramble ablation} The prompt is left
byte-identical except that the drug name and mechanism are replaced by another
drug's, the control cell and the scoring truth unchanged
(\cref{fig:scramble}). Control- and batch-derived leakage therefore cancels in
the scramble gap $\gap = \NIR_{\textrm{model}} - \NIR_{\textrm{scramble}}$, the
only difference in this chapter whose leakage cancels by construction. It does
not displace the drug-agnostic comparisons above and below, which rest on the
within-plate rule instead.

% PLACEMENT ONLY. The float was anchored six paragraphs lower (just before "The
% third instrument"), which is where its [htbp] first became satisfiable -- the
% top of the NEXT page -- and that page then had no room left for
% \cref{tab:blind}, which was pushed a further page past the sentence that
% introduces it (measured: anchor p35, caption p37). Anchoring the figure at the
% paragraph it actually depicts, the scramble ablation itself, is earlier in the
% flow, so the figure takes the top slot one page sooner and the table follows on
% the page after its own anchor. Nothing about the figure or its caption changes.
%
% The paragraph above now carries the \cref, so the anchor is explicit and not
% merely positional. Its second sentence lost the restatement "Both arms share
% the same control cell, so ..." because the figure now sets that clause on one
% unbroken line beneath the two arms; the inference is carried by "therefore"
% off "the control cell ... unchanged" in the sentence before it. No claim,
% hedge or magnitude changed, and the paragraph is one line shorter than it was.
%
% \input, not \includegraphics: figs/fig-scramble.tex is TikZ and carries its
% own float, caption and \label. \input resolves against the main-document
% directory, so this is thesis_v2/figs/, NOT the read-only ../thesis/figs/ that
% \graphicspath adds for raster/PDF assets only.
% ===========================================================================
% fig-scramble.tex (v2) -- the scramble ablation and why its difference is
% leak-immune. NO DATA DEPENDENCY: this is a schematic. The two drug names and
% the two mechanism strings are the illustrative pair already used in v1; the
% only quantity in the float is the definition
% \gap = \NIR_real - \NIR_scrambled, which is the definition given in
% Sections/04a-ruler.tex, not a measurement.
%
% CHANGES FROM v1 (thesis/figs/fig-scramble.tex), all layout, none semantic:
%   (1) the closing sentence is now ONE unbroken line. In v1 it was set as two
%       centred nodes with a dashed arrow BETWEEN them, so the reader crossed a
%       rule in the middle of the clause "so control- and batch-derived /
%       leakage cancels". Measured at \scriptsize the sentence is 360.25pt,
%       which fits the 404.03pt measure with 43.8pt to spare, so the break was
%       never forced -- it existed only to make room for the arrow.
%   (2) the dashed arrow is deleted. It ran left-to-right under the figure and
%       connected nothing: no node at either end, no quantity implied by its
%       direction.
%   (3) the \gap definition is no longer floating in the white between the two
%       curved connectors. It hangs off the LOWER connector on a short vertical
%       tick, so the eye reads "these two arms, differenced" rather than
%       finding a formula parked in open space.
%
% WIDTH. v1 measured 289.4pt = 101.7mm inside a 142mm measure and sat visibly
% short of the block. This version is laid out to 14.2cm = 142.0mm = the v2
% \textwidth (404.03 TeX pt, preamble.tex:88). It is NOT \resizebox'd and it is
% NOT [scale=]'d: every type size below is the size that prints. Widening was
% done by pinning the truth node to the right edge of the block (anchor=east at
% x=14.2) and widening the prompt cards from their natural 3.70cm to 5.60cm, so
% the extra 40mm goes into the converging arms, which is the part of the
% picture that carries the argument, and not into the card interiors.
%
% MEASURED (\settowidth, 12pt Pagella, this preamble):
%   \textwidth                                     404.03pt
%   closing sentence, \scriptsize                  360.25pt   -> one line, fits
%   closing sentence, \footnotesize                450.31pt   -> would not fit
%   prompt tabular, natural                        104.88pt
%   "scrambled prompt --- only the two ..."        192.22pt = 6.78cm
%   \gap definition, \small                        137.96pt
% Those last two set the card width. The cards stop at 5.60cm and the truth
% node starts at 12.09cm, leaving 6.49cm of gap; the 4.86cm \gap definition
% hangs in the middle of that gap with 0.83cm clear on its left and 0.68cm on
% its right, so it touches neither the card nor the truth. A wider card (6.90cm,
% flush with the 6.78cm label above it) was tried and rejected: it squeezed
% those clearances to 3pt and 2pt.
% ===========================================================================
\begin{figure}[htbp]
\centering
\begin{tikzpicture}[
  x=1cm, y=1cm, font=\small,
  lbl/.style={font=\scriptsize, text=quietink},
  strong/.style={font=\scriptsize, text=ink},
  % text width, not minimum width: minimum width would centre the tabular in
  % the enlarged box and the prompt would stop looking like a prompt.
  promptbox/.style={draw=rulegrey, rounded corners=2pt, inner sep=5pt,
                    fill=rulegrey!8, text width=5.248cm, align=left},
  truthbox/.style={draw=rulegrey, rounded corners=2pt, inner sep=6pt,
                   align=center, fill=rulegrey!10},
  ar/.style={-{Stealth[length=4.5pt]}, quietink},
  tick/.style={quietink, line width=0.5pt},
]
% rulegrey!35 reproduces v1's black!14 span shading against the page's own
% rule ink rather than pure black.
\newcommand{\spanhl}[1]{\colorbox{rulegrey!35}{#1}}

% ---- arm 1: the real prompt --------------------------------------------
\node[lbl, anchor=west] at (0,0) {real prompt};
\node[promptbox, anchor=north west] (p1) at (0,-0.26) {%
  \renewcommand{\arraystretch}{0.9}%
  \begin{tabular}{@{}l@{~~}l@{}}
  \ttfamily\scriptsize Cell line:    & \ttfamily\scriptsize MCF7 \\
  \ttfamily\scriptsize Drug:         & \ttfamily\scriptsize \spanhl{Lapatinib} \\
  \ttfamily\scriptsize Mechanism:    & \ttfamily\scriptsize \spanhl{EGFR inhibitor} \\
  \ttfamily\scriptsize Control cell: & \ttfamily\scriptsize ACTB GAPDH \ldots \\
  \end{tabular}};

% ---- arm 2: the scramble. Positioned OFF p1's measured south edge, never off
% an estimate of the card height -- the card is ~2.2cm tall once the shaded
% spans are set, which is more than four lines of \scriptsize looks like it
% should need.
\node[lbl, anchor=north west] (l2) at ($(p1.south west)+(0,-0.34)$)
  {scrambled prompt --- only the two shaded spans differ};
\node[promptbox, anchor=north west] (p2) at ($(l2.south west)+(0,-0.10)$) {%
  \renewcommand{\arraystretch}{0.9}%
  \begin{tabular}{@{}l@{~~}l@{}}
  \ttfamily\scriptsize Cell line:    & \ttfamily\scriptsize MCF7 \\
  \ttfamily\scriptsize Drug:         & \ttfamily\scriptsize \spanhl{Vorinostat} \\
  \ttfamily\scriptsize Mechanism:    & \ttfamily\scriptsize \spanhl{HDAC inhibitor} \\
  \ttfamily\scriptsize Control cell: & \ttfamily\scriptsize ACTB GAPDH \ldots \\
  \end{tabular}};

% ---- the single truth, on the right edge of the measure -----------------
% anchor=east at x=14.2 pins the picture's right edge to \textwidth exactly.
\coordinate (mid) at ($(p1.east)!0.5!(p2.east)$);
\node[truthbox, anchor=east] (truth) at (14.2,0 |- mid)
  {truth for\\\textbf{Lapatinib}\\{\scriptsize drawn once}};

% Both arms leave horizontally and arrive on the truth from the outside, so
% the convergence is the shape of the picture.
\draw[ar] (p1.east)
  .. controls ($(p1.east)+(2.7,0)$) and ($(truth.north west)+(-1.0,0.30)$)
  .. (truth.north west);
\draw[ar] (p2.east)
  .. controls ($(p2.east)+(2.7,0)$) and ($(truth.south west)+(-1.0,-0.30)$)
  .. (truth.south west) coordinate[pos=0.42] (hang);

% ---- the definition, hung off the lower connector ----------------------
\draw[tick] (hang) -- ($(hang)+(0,-0.62)$);
\node[anchor=north, font=\small, inner sep=1pt] (gapdef) at ($(hang)+(0,-0.62)$)
  {$\gap = \NIR_{\text{real}} - \NIR_{\text{scrambled}}$};

% ---- the reason the difference is leak-immune, on one line -------------
\node[strong, anchor=north] at (7.1,-6.00)
  {the shared control cell enters \emph{both} arms, so control- and
   batch-derived leakage cancels in the difference};

\end{tikzpicture}
\caption[The scramble ablation]{\textbf{The scramble replaces the drug name and its mechanism and
nothing else.} Both arms are conditioned on the same control cell and scored against the same truth,
so any leakage through the control cell or the plate enters both arms in equal measure and cancels in
the difference $\gap$. This is why $\gap$ is the only difference in the chapter whose leakage cancels
by construction, while a comparison against a drug-agnostic predictor is not leak-immune --- such a
baseline inherits whatever plate structure the comparison set carries
(Table~\ref{tab:plateleak}). Comparisons of that kind are used, and used centrally, with the
comparison set declared at every score: within plate in the expression frame, where the leakage is
controlled rather than cancelled, and by cell line in the residual frame, where the plate-scoped
generic has already subtracted each plate's mean response from truths and predictions alike --- in
expectation, with however much survives in an individual residual left unmeasured
(\S\ref{sec:methods-residual}, and \S\ref{sec:limitations} for what that leaves open).
A refinement matters here: if the swapped-in drug happens to have a
\emph{similar} true response then an unchanged output is correct rather than blind, so the swap
partner is chosen from three defined similarity strata by residual cosine. Which stratum can carry
the canonical magnitude associated with this figure is constrained by the comparator's own score.
\S\ref{sec:res-generalisation} finds the geometry-defined \texttt{orth} partner empirically near
chance at $0.501$ in this sample, making it the least contaminated available comparator rather than
an independently established null. That restriction belongs to this canonical scramble comparison:
the three-stratum repair comparison differences each model score against its own
\texttt{near}/\texttt{orth}/\texttt{opposite} partner, while field attribution uses its stated
span-specific partner and retains the resulting comparator caveat.}
\label{fig:scramble}
\end{figure}


% fig:comparator sits here, two paragraphs above the sentence that cites it,
% because both floats then queue from the same point and LaTeX keeps them in
% order: fig:scramble takes the next top slot and fig:comparator the one after,
% which lands it on the spread that carries the stratified-scramble paragraph.
% It is cited here and used in full six sections later (sec:res-generalisation);
% placing it there instead would put the picture of what the strata are three
% chapters' worth of pages after the paragraph that defines them. Its caption
% carries the forward pointer, so the reader who meets it early is told where
% the argument is finished. Same \input rule as above: figs/fig-comparator.tex
% carries the float, caption and \label, and reads fig-comparator.pdf out of
% thesis_v2/figs/.
% fig-comparator.tex -- float wrapper for fig:comparator.
% Data: thesis_v2/figs/fig-comparator.json, written by thesis_v2/figs/comparator.py from
% RESULTS_cluster/scramble_stratum_audit_v4.json and RESULTS_cluster/re_v3.json.
% Placement: Sections/04a-ruler.tex, immediately after \input of fig-scramble.
% NO width= KEY. The PDF is drawn at 142mm, which is exactly \linewidth in the body block, so
% \includegraphics{} places it at 1.0 and the 10pt type in it prints at 10pt. Adding
% [width=\linewidth] would be harmless today and a silent 11% type shrink the moment the measure
% changes; \resizebox would be wrong at any measure.
\begin{figure}[htbp]
\centering
\includegraphics{fig-comparator}
% Short caption: 46 characters, and it states the finding rather than the topic. The long form
% below carries the calibration half of the argument; the List of Figures gets the consequence.
\caption[Most of the \texttt{opposite} gap is the comparator]{%
\textbf{The comparator's own score is a monotone function of how correct the swapped-in partner
is, so most of the \texttt{opposite} stratum's headline gap is the comparator failing rather than
the model knowing.} All quantities are residual-frame $\NIR$ scored against the cell-line
comparison set, pooled over the $\num{1394}$ scored conditions.
\emph{(a)} Each mark is one swap stratum, at its mean partner residual cosine and its mean
$\NIR$. The score falls from $0.5504$ through $0.5006$ to $0.4228$ as the partner's true residual
moves from aligned to anti-aligned, so \texttt{near} hands the model a partly correct answer and
\texttt{opposite} an actively wrong one. The straight line is the chord joining the two extreme
marks, drawn as a guide to the eye; it is not a fit, and \texttt{orth} lies $0.016$ above it.
Because the strata are defined using the \emph{true} residual geometry, \texttt{orth}'s sitting on
chance is an empirical property of this sample and not an independently established null: it is
the least contaminated comparator available, which is a weaker statement
(\cref{sec:res-generalisation}).
\emph{(b)} The same gap $\gap$ split at chance into its two terms,
$\gap = (\NIR_{\textrm{model}} - 0.5) + (0.5 - \NIR_{\textrm{scramble}})$. The accent bar is the
model's own advantage over chance, $+0.0367$, and is identical on all three rows because it is the
same model arm; the strata differ only in the grey bar, the comparator's own displacement from
chance, drawn on the side it actually falls --- right of the rule it subtracts from the gap, left
of the rule it adds. On \texttt{opposite} the grey term is $+0.0772$ of the $+0.1139$ reported,
$68\%$ of it, which is the annotation on that row.
Intervals are $95\%$ two-way cluster-robust on cell line and treatment well ($42$ cell lines,
$200$ wells); the one-way variants stored alongside them are not drawn.}
\label{fig:comparator}
\end{figure}


\paragraph{Two constraints on the swap partner} The partner must be a different
drug. Allowing the same drug at another dose or plate substitutes the drug for
itself, an unchanged prediction is then correct and not blind, and the contamination
concentrates in the most similar stratum. It must also be on the same plate,
since residuals retain plate structure and a cross-plate partner would differ in
batch as well as in drug.

A scramble that swaps in a drug with a similar true response is a weak test, and
that risk is structural here. In full-profile space the mechanism annotation is
close to uninformative about transcriptional response. Same-mechanism drugs in
Tahoe are just over two per cent closer to each other than arbitrary pairs
are.\seealso{\cref{sec:res-encoding} quotes the distances behind that statement}
Hence conditioning on mechanism fails to move $\DEdr$ in the grouping ladder of
\cref{sec:res-lookup}, and a partner of different mechanism cannot be relied on
for a different response. Mechanism does carry signal in the residual frame
(\cref{sec:res-scope}); in full-profile space the generic response swamps it, a
frame difference and not a contradiction.

Two remedies are used. A swap-distance sweep bins (drug, swap) pairs by the
response distance $\lVert y_A - y_B \rVert$ between real and swapped-in drug,
turning the single test into a curve (\cref{sec:app-sweep}). And a stratified
scramble defines three partners per condition by residual cosine within the cell
line: \texttt{near} (most similar), \texttt{orth} ($\cos\approx0$) and
\texttt{opposite} (most anti-correlated). The strata expose what a single
comparator hides (\cref{fig:comparator}): the gap it defines has two parts, how
much naming the truth helps and how much naming a lie hurts, and only the first
is wanted, so a usable null is a stratum scoring at chance.
\Cref{sec:res-generalisation} shows exactly one qualifies, and not the one an
evaluation would default to.

\paragraph{The third instrument is forced-choice grading} A prediction is scored
against its own condition's truth and against another drug's truth from the same
comparison set, recording how often its own wins. Chance is $0.50$ by
construction, and the precision reference lands between $0.67$ and $0.83$
depending on cells per condition.

\paragraph{The fourth is output invariance} It asks whether the drug token
changes the generated text at all, without reference to any truth. It is run at
$3800$ tokens, the longest budget in the chapter, because
truncation would cut the comparison
itself.\seealso{\cref{sec:methods-data} records the four generation budgets} For
a fixed control cell it compares
% \resizebox, not a reflow: the two \underbrace labels are the width, and both
% are wanted on one line so the subtraction reads as one sentence. Natural
% width is 165.6mm against the 142mm block (67.15pt over), so this sets at
% about 0.86. Nothing here is body text.
\begin{equation}
\resizebox{\linewidth}{!}{$\displaystyle
 \textrm{invariance gap} \;=\;
 \underbrace{\overline{\textrm{sim}}(\textrm{same prompt, resampled})}%
   _{\textrm{sampling noise alone}}
 \;-\; \underbrace{\overline{\textrm{sim}}(\textrm{prompts differing in the
   drug})}_{\textrm{noise plus any drug effect}},
$}
\end{equation}
with similarity as top-$100$ Kendall $\tau$ and as Jaccard overlap, over $8$
contexts $\times$ $12$ drugs at temperature $0.8$. It is the least informative
of the four and the least confoundable.

\begin{table}[htbp]
\begin{fullwidth}
\begin{threeparttable}
\caption{\textbf{Four independent instruments agree, and the two that return
intervals bound the drug's effect instead of merely covering zero.} Near-ceiling
$\DEdr$ is achieved identically with a wrong-mechanism drug, so that score is
drug-independent.}
\label{tab:blind}
\small
% Column 3 was `l'. At the widened measure its natural width ("$0.48$
% ($\approx$ scramble; reference 0.67--0.83)") plus columns 1--2 left the X
% column too narrow for the single unbreakable word "indistinguishable", which
% ran 43.91pt past the fullwidth block. Column 3 is now a weighted X so it
% wraps; the two \hsize factors sum to 2, which is what tabularx requires of
% two X columns.
\begin{tabularx}{\linewidth}{@{}l l
  >{\hsize=1.20\hsize\raggedright\arraybackslash}X
  >{\hsize=0.80\hsize\raggedright\arraybackslash}X@{}}
\toprule
\thead{Instrument} & \thead{Quantity} & \thead{Value} & \thead{Reading} \\
\midrule
Forced-choice grading & own-truth win rate
  & $0.48$ ($\approx$ scramble; reference 0.67--0.83) & chance \\
Output invariance & invariance gap & \ci{0.000}{-0.019}{+0.016}
  & bounded, $\leq 0.016$ \\
\addlinespace
Scramble, tier 1 & $\DEdr$ & $0.740$ real vs $0.739$ scrambled
  & indistinguishable \\
Scramble, tier 2 & $\DEdr$ & $0.733$ real vs $0.729$ scrambled
  & indistinguishable \\
Scramble (identifiable subset) & $\gap$ ($\NIR$) & \ci{+0.014}{-0.016}{+0.042}
  & bounded, $\leq +0.042$ \\
\bottomrule
\end{tabularx}
\begin{tablenotes}[flushleft]\footnotesize
\item The two $\DEdr$ rows are paired comparisons on one scale, the only use
\cref{sec:res-metrics} licenses for that metric.
\end{tablenotes}
\end{threeparttable}
\end{fullwidth}
\end{table}

% PLACEMENT ONLY. This is one of the three survivors the \paragraph orphan patch
% in preamble.tex section 4 records: the global reservation is
% 3\baselineskip, which this head cleared with 2.3 lines to spare and then set
% itself plus exactly two lines at the foot of the page. 5 lines is head plus
% four, which the page cannot fake.
\needspace{5\baselineskip}%
\paragraph{An interval that covers zero earns its keep by what it excludes} Not
by the word \emph{null}. Two predictions from a byte-identical prompt agree at
$\tau = 0.198$, so the invariance interval caps any drug-driven divergence at
roughly a twelfth of that agreement, and the Jaccard version is tighter still,
\ci{+0.000}{-0.008}{+0.008}. On the identifiable stratum the scramble gap of
\ci{+0.014}{-0.016}{+0.042} runs against the $0.170$ separating the model's
$0.768$ from that stratum's $0.938$ reference; at the very top of the interval,
naming the right drug would buy a quarter of what is there to be bought, and the
point estimate puts it at $8\%$ of that. Neither figure shows the effect is
zero, and no experiment here could. Both put it far below the headroom these
conditions leave.

\paragraph{Four objections close on measurement, not argument} The first is
whether there is anything to know at all. A ladder of mean-shift baselines
answers it, and the informative conditioning variable is the cellular context,
not the drug class.\seealso{\cref{sec:res-lookup} gives the ladder}

The second decides whether this section reports a result or an instrument
failure. Barely half the atlas clears the pre-set identifiability bar
(\cref{tab:scale}), so the verdict is re-scored where the wells are demonstrably
adequate. That does not rescue the model. It appears to beat the linear baseline
there ($0.768$ against $0.531$), yet a zero-drug-information control-copy
reaches $0.766$ on the same conditions, the paired contrast is
\ci{+0.002}{-0.026}{+0.028}, and the leak-immune $\gap$ stays bounded as above.

The control arm carries a second result that is easy to walk past. Across the
three strata its score tracks the reference almost step for step
(\cref{fig:difficulty}, panel b): $0.310$
against a reference of $0.287$ on the low-reference stratum, $0.540$ against
$0.689$ on the marginal one, $0.766$ against $0.938$ on the identifiable one. A
predictor knowing nothing whatever about any drug does best exactly where the
ranking problem is easiest, so what makes a condition identifiable is largely
what makes its untreated cells distinctive among their plate-mates, and the
distance to a high reference is not all room that better drug modelling would
fill. The reference rises with aggregation (\cref{fig:difficulty}, panel a:
$0.61$ at $10$ cells, $0.76$
at $40$, $0.85$ at $120$); what the strata add is that most of the ranking can
be done without the drug.

\begin{figure}[htbp]
\centering
% drawn at 142mm by thesis_v2/figs/difficulty.py; no width= key, so it is placed
% at scale 1.0 and its type prints at the point sizes it was set in.
\includegraphics{fig-difficulty}
\caption[Difficulty structure and the stratum ladder]{\textbf{Conditions are
more empirically discriminable at the tested pseudobulk budget where the
precision reference is high, and precisely there, a predictor containing no drug
information matches the model.} Both panels are $\NIR$ in the expression frame,
on one shared axis. (a) The within-well split-half precision reference against
cells per condition, for two comparison sets: cross-plate, and within plate. (b)
The same reference used to bin conditions, five arms per stratum, $n = 296$
unwinnable, $106$ marginal, $204$ identifiable. On the unwinnable stratum the
reference itself is $0.287$, below chance, as are the model, control-copy and
scramble marks. On the identifiable stratum the model reaches $0.768$ and a
control-copy baseline $0.766$, paired contrast \ci{+0.002}{-0.026}{+0.028},
clustered over cell lines. Plotting the model alone would suggest it does better
on higher-reference drugs. The arm marks are bare: the source carries no
interval for a stratum mean, and this document draws only the clustered
bootstrap, so nothing is drawn rather than an interval of another kind.}
\label{fig:difficulty}
\end{figure}

The third objection is whether the representation throws the drug away. Cell
sentences keep only the rank order of genes and discard expression magnitude; if
the drug lived there the following sections would attack the wrong thing. It is
testable without any model, by asking whether true expression discriminates
drugs better than rank does, and it does not.\seealso{\cref{sec:res-lookup}
gives that measurement}

The fourth is whether the scramble swaps in similar drugs. The swap-distance
sweep bounds this rather than settling it, and it sits outside the within-plate
rule, being an earlier run at \texttt{max\_new\_tokens=1200} scored on the
unseen-drug tier of a different evaluation set, with swap partners not
plate-restricted, and with $n=40$ per bin wide enough to exclude a large effect
but not to establish a null (\cref{sec:app-sweep}). Neither column has an
interval excluding zero in any distance bin; in the single-cell column the far
tail reads $-0.019$ across a $2.3\times$ spread in swap distance, trend
correlation $+0.051$. That column bounds a large dissimilarity effect for the
expression-frame model. The optimal-transport column cannot, since
\cref{sec:res-encoding} later shows that checkpoint clears zero under the
canonical within-plate stratified test at the matched
\texttt{max\_new\_tokens=1400} budget. Its apparent null is budget-dependent,
and the stratified scramble replaced the sweep in the reported set.

\begin{claim}
The model is at chance overall ($0.498$ against a $0.576$ within-well split-half
precision reference), matched on the identifiable stratum by a
zero-drug-information control-copy ($0.768$ against $0.766$, paired contrast
\ci{+0.002}{-0.026}{+0.028}), and no more moved by a maximally dissimilar swap
than by a similar one in the single-cell column of a non-canonical
\texttt{max\_new\_tokens=1200} sweep outside the within-plate rule
(\cref{sec:app-sweep}). The sweep's optimal-transport column does not carry that
bound; \cref{sec:res-encoding} later shows the same checkpoint is non-null under
the canonical test at the matched $1400$-token budget.
\end{claim}
